\author{Tim Feldmüller}
\title{Das Framing von Extremismusvarianten im medialen Diskurs der Jahre 1999–2021}
\subtitle{Eine \textit{corpus-driven} Methode zur Erschließung und Visualisierung semantischer Frames}
\BackBody{Die vorliegende Arbeit untersucht das Framing von Extremismusvarianten im medialen Diskurs der Jahre 1999–2021 und entwickelt eine \emph{corpus-driven} Methode zur Erschließung und Visualisierung semantischer Frames. Sie verbindet Frame-semantische Theorie nach Busse mit Methoden und Annahmen der distributionellen Semantik und entwickelt ein Verfahren, um diskursives Wissen aus der musterhaften Verteilung von Wörtern in Korpora induktiv zu rekonstruieren.
Als Datengrundlage dient ein umfangreiches Korpus aus Online-Artikeln der Zeitungen \textsc{taz}, \textsc{Spiegel} und \textsc{Welt}. Methodisch kombiniert die Arbeit Word Embeddings und Kollokationsanalysen, um Frame-Elemente und ihre Relationen zu identifizieren und als Kollokationsnetzwerke zu visualisieren. Dieser Ansatz ermöglicht es, einen großen Teil des verstehensrelevanten Wissens zu erschließen, das zum Verständnis von Extremismus als sozialem Konzept notwendig ist, ohne auf vorab definierte Kategorien zurückzugreifen.
Die Analyse untersucht kontrastiv das Framing von Rechts- und Linksextremismus sowie Islamismus über acht Zeiträume hinweg sowie im Vergleich der Zeitungen. Sie zeigt, wie unterschiedliche Extremismusvarianten in den drei Zeitungen gerahmt werden, welche Akteure als extremistisch gelten, welche Handlungen ihnen zugeschrieben werden und wie sich diese diskursiven Konstruktionen im Zeitverlauf wandeln. Die Arbeit leistet damit sowohl einen methodischen Beitrag zur korpuslinguistischen Diskursanalyse als auch einen inhaltlichen zur Beschreibung des medialen Extremismusdiskurses in Deutschland.
}
\ISBNdigital{978-3-96110-557-1}
\ISBNhardcover{978-3-98554-177-5}
\BookDOI{10.5281/zenodo.18274343}
\typesetter{Tim Feldmüller, Hannah Schleupner}
\proofreader{Jean Nitzke, Ludger Paschen, Niklas Reinken, Hannah Schleupner, Rainer Schulze}
\lsCoverTitleSizes{40pt}{13mm}
\renewcommand{\lsSeries}{frats}  
\renewcommand{\lsSeriesNumber}{2}
\renewcommand{\lsID}{567}
